% Chapter 12: The Geometric Langlands Program
\let\LocalMacro\relax

\chapter{Geometric Langlands Program}

\section{The Problem}

\subsection{Informal}
Mathematics has two huge areas that seem to have nothing in common: one studies numbers and equations, the other studies shapes and curves. Over the past century, mathematicians began to suspect that there might be a perfect hidden dictionary between these two areas---a way to translate every question in one into an equivalent question in the other. Proving that this dictionary exists has become one of the deepest and most active quests in modern mathematics.

\subsection{Formal}
The geometric Langlands conjecture posits a correspondence between:
\begin{itemize}
\item \emph{Automorphic objects}: algebraic sheaves on the stack $\mathrm{Bun}_G$ of $G$-bundles on a smooth projective curve $X$
\item \emph{Galois objects}: D-modules (or local systems) on the stack $\mathrm{Bun}_{^L G}$ of $^L G$-bundles on $X$
\end{itemize}
where $^L G$ is the Langlands dual group of $G$. The conjecture asserts that these two categories are equivalent as $\ell$-adic derived categories.

\subsection{Historical Context}
The geometric Langlands program was proposed by Drinfeld in 1987 and developed by Beilinson and Drinfeld in the 1990s \cite{drinfeld1987, beilinson2004}. It has become one of the deepest and most active areas of mathematics, connecting algebraic geometry, representation theory, and mathematical physics (particularly topological field theory).

\subsection{What Was Known}
The classical Langlands program, proposed in the 1960s, conjectured a correspondence between number theory and the theory of automorphic forms. The geometric version, proposed by Drinfeld in 1987, reinterpreted this correspondence in the language of geometry. By the 2010s, partial results had been established for specific groups and specific cases, but the full conjecture remained wide open.

\subsection{Why It Matters}
A perfect dictionary between number theory and geometry would mean that every hard problem in one area could be translated into the other, where different tools might crack it. The geometric Langlands program has already led to major advances in the understanding of fundamental forces in physics and the deep structure of mathematical truth.

\subsection{Simplified Version}
The classical (arithmetic) Langlands program, proposed in the 1960s, seeks a correspondence between automorphic forms on $GL_n$ over $\mathbb{Q}$ and $n$-dimensional Galois representations. The geometric version replaces number fields with function fields and replaces individual forms with sheaves (geometric objects spread over a space). The simplest case---the correspondence for $GL_1$ over a curve---was already understood classically; the deep results of the geometric program concern $GL_n$ for $n \ge 2$.

\section{Mathematical Theory}


\subsection{Prerequisites}
This chapter assumes familiarity with:
\begin{itemize}
\item Algebraic geometry (sheaves, schemes, stacks)
\item Representation theory (D-modules, categories of sheaves)
\item Category theory (derived categories, equivalences)
\end{itemize}

\subsection{The Stack of Bundles}

\begin{definition}[Stack of $G$-Bundles]
Let $X$ be a smooth projective curve over a field $k$ of characteristic 0, and let $G$ be a reductive group. The stack $\text{Bun}_G(X)$ classifies $G$-bundles on $X$.
\end{definition}

The geometric Langlands program studies the derived category $D(\text{Bun}_G(X))$ of D-modules (or constructible sheaves in characteristic $p$) on this stack.

\subsection{Local Systems}

\begin{definition}[Local System]
A \emph{local system} on $X$ with fiber in $^LG$ (the Langlands dual group) is a flat vector bundle, equivalently a representation of the fundamental group $\pi_1(X)$.
\end{definition}

The stack of local systems is denoted $\text{LocSys}_{^LG}(X)$.

\subsection{The Duality Conjecture}

\begin{conjecture}[Geometric Langlands Duality]
There is an equivalence of categories:
\begin{equation}
\mathbb{L}: D(\text{Bun}_G(X)) \xrightarrow{\sim} D(\text{LocSys}_{^LG}(X))
\end{equation}
where the left side is the derived category of D-modules on the stack of $G$-bundles, and the right side is the derived category of coherent sheaves on the stack of local systems for the Langlands dual group $^LG$.
\end{conjecture}

This is a categorical generalization of the classical Langlands correspondence.

\subsection{What Are These Objects?}
Before diving into the proof, here is a brief guide to the objects appearing in the conjecture:

\begin{itemize}
\item \textbf{D-modules}: Think of a D-module as a ``differential equation on a space.'' Just as a function on $\mathbb{R}$ can satisfy $f'(x) = f(x)$, a D-module on a geometric space $X$ encodes a system of differential equations that the ``sections'' must satisfy. D-modules are the geometric analogue of representations of Galois groups.
\item \textbf{A $G$-bundle}: A $G$-bundle on a curve $X$ is a space that looks locally like $G \times X$ (a product of the group $G$ with the curve), but may be ``twisted'' globally. Think of it as assigning a copy of the symmetry group $G$ to each point of the curve, with a rule for how to compare copies at nearby points. The stack $\mathrm{Bun}_G$ parameterizes all such bundles.
\item \textbf{Local systems}: A local system on $X$ is a way of assigning a vector space to each point of $X$ and transporting vectors along paths---like a ``parallel parking'' rule for vectors. The stack $\mathrm{LocSys}_{^LG}$ parameterizes all local systems for the dual group ${}^LG$.
\item \textbf{Derived categories}: Instead of studying individual D-modules or sheaves, we study their \emph{relationships} (morphisms). The derived category records these relationships up to homotopy, keeping track of higher information that ordinary categories discard.
\item \textbf{The Langlands dual group ${}^LG$}: For each reductive group $G$ (like $\mathrm{GL}_n$), there is a dual group ${}^LG$ constructed from the root data of $G$. For $\mathrm{GL}_n$, the dual group is again $\mathrm{GL}_n$. For other groups, the dual can be different (e.g., $\mathrm{SL}_n$ is dual to $\mathrm{PSL}_n$).
\end{itemize}

In plain language: the conjecture says that the space of ``differential equations on $G$-bundles'' is secretly the same as the space of ``parallel-transport rules for ${}^LG$-representations.'' The two descriptions look completely different, but the conjecture (now theorem) says they carry exactly the same information.

\section{The Solution}

\subsection{The Big Idea}
The geometric Langlands conjecture says: the derived category of D-modules on $\mathrm{Bun}_G$ is equivalent to the derived category of coherent sheaves on $\mathrm{LocSys}_{^LG}$. In plain terms (see ``What Are These Objects?'' above): the space of differential equations on $G$-bundles is the same as the space of parallel-transport rules for ${}^LG$-representations. The 2024 proof builds this equivalence in two steps. First, the \textbf{geometric Satake equivalence} (Mirkovi\'c--Vilonen) provides a local dictionary at each point of the curve: it identifies representations of ${}^LG$ with certain sheaves on the affine Grassmannian of $G$. Second, \textbf{factorization categories} glue these local dictionaries into a global equivalence. The proof required over 1,000 pages because the ``gluing grammar''---the categorical framework of ind-coherent sheaves on stacks---is vastly more sophisticated than ordinary category theory.

\subsection{What Was Stopping Progress}
The classical Langlands program had established a dictionary between number-theoretic and geometric objects for the simplest case ($GL_1$) and some special groups, but the full conjecture---a categorical equivalence for \emph{all} reductive groups---required mathematical machinery that simply did not exist. The right language was not ordinary categories but \emph{derived categories of ind-coherent sheaves on stacks}, a vastly more sophisticated framework capable of handling singularities, non-properness, and infinite-dimensional objects. Previous approaches by Beilinson, Drinfeld, Laumon, and Rothstein built important pieces of the dictionary but could not assemble them into a complete translation. It was as if linguists had decoded individual words and grammar rules in two languages but had never produced a complete grammar that could translate any sentence.

\subsection{The Breakthrough}
Gaitsgory's program, culminating in the 2024 collaboration of nine mathematicians, introduced \textbf{ind-coherent sheaves} as the correct categorical framework and proved the equivalence via the \textbf{spectral gluing construction}. The key idea: decompose the Langlands duality into local pieces using the \emph{geometric Satake equivalence} (which handles the ``dictionary entry'' at each point of the curve) and then glue these pieces together using the theory of \emph{factorization categories}. The proof required over 1,000 pages across five preprints, combining techniques from derived algebraic geometry, topological field theory (Chern--Simons duality), and higher category theory. The result: an equivalence of derived categories between D-modules on the stack of $G$-bundles and coherent sheaves on the stack of ${}^L\!G$-local systems, for any reductive group $G$.

\subsection{How It Works}
The proof proceeds in two stages. First, the \textbf{geometric Satake equivalence} (proved by Mirkovi\'c--Vilonen) provides a local dictionary at each point of the curve $X$: it identifies the derived category of $G$-equivariant D-modules on the affine Grassmannian of $G$ with the derived category of representations of the Langlands dual group ${}^L\!G$. Intuitively, at each point of the curve, there is a perfect local translation between $G$-bundles and ${}^L\!G$-representations. Second, the theory of \textbf{factorization categories} (or factorization algebras) glues these local translations into a single global equivalence. Think of having a word-by-word dictionary for every position in a sentence: the factorization structure is the grammar that assembles these local words into a coherent global meaning. The technical heart of the proof is showing that the global functor built from local Satake data is independent of choices---this requires the language of ind-coherent sheaves on stacks, which handles singularities and infinite-dimensional objects that ordinary categories cannot.

\subsection{What Was Missing}
The classical Langlands program posits a correspondence between automorphic forms and Galois representations for reductive groups over number fields. The geometric version, proposed by Drinfeld in 1987, reinterpreted this in the language of algebraic geometry: an equivalence between D-modules on the stack of $G$-bundles and coherent sheaves on the stack of ${}^L\!G$-local systems. What was missing was a proof that this equivalence holds for \emph{all} reductive groups---previous results covered $GL_1$ and some special cases, but the full conjecture required categorical machinery that did not yet exist.

\subsection{Why Existing Methods Failed}
The classical techniques of Beilinson--Drinfeld established the framework for $GL_1$ and identified the Hecke eigensheaf condition as the geometric analogue of an automorphic form, but their methods relied on ad hoc constructions that did not generalize. Laumon and Rothstein extended the theory using the geometric Satake equivalence, but only for specific groups and local settings. The fundamental obstruction was that the correct \emph{categorical} language---derived categories of ind-coherent sheaves on stacks---was not sufficiently developed to formulate, let alone prove, the full duality.

\subsection{What Was Novel}
Gaitsgory's program introduced \textbf{ind-coherent sheaves} as the correct categorical framework, replacing ordinary D-modules with a richer theory that handles singularities and non-properness. The 2024 proof by Arinkin, Gaitsgory, Raskin et al.\ (five preprints, $1000+$ pages) established the unramified geometric Langlands conjecture for any reductive group $G$ using the \textbf{spectral gluing construction}: the Langlands dual is realized by gluing local data via the geometric Satake equivalence (proved by Mirkovi\'c--Vilonen) and the theory of factorization categories. Higher category theory ($\infty$-categories) and topological field theory (Chern--Simons duality) provide the structural backbone.

\subsection{Student-Friendly Explanation}
The Langlands program says there is a dictionary between two apparently unrelated mathematical worlds: the world of symmetries (automorphic forms) and the world of algebraic equations (Galois representations). In the geometric version, both sides become geometric objects spread over a space. The conjecture says these two collections of geometric objects are ``the same''---you can translate back and forth without losing information. The proof required building a vast categorical machine: think of it as constructing a bilingual library where every book on one side has an exact counterpart on the other, with the cataloging system itself (the category theory) ensuring nothing is lost in translation.

\subsection{Beilinson--Drinfeld (1990s)}

Beilinson and Drinfeld laid the foundations of the geometric Langlands program, establishing \cite{beilinson2004}:

\begin{itemize}
\item The correct categorical framework (D-modules on stacks)
\item The role of the Hecke eigensheaf condition (geometric analogue of automorphic forms)
\item The connection to integrable systems (the Hitchin fibration)
\end{itemize}

They proved the correspondence for $GL_1$ and established the framework for $GL_n$.

\subsection{Ben-Zvi, Francis, Nadler (2010s)}

David Ben-Zvi, John Francis, and David Nadler developed a comprehensive framework for the geometric Langlands program using the language of \emph{higher categories} and \emph{topological field theory} \cite{benzvi2014}:

\begin{theorem}[Ben-Zvi, Francis, Nadler, 2010s \cite{benzvi2014}]
The geometric Langlands correspondence can be understood as a manifestation of the duality between 3-dimensional topological field theories (specifically, the Chern--Simons theory and its dual).
\end{theorem}

Their approach:

\begin{enumerate}
\item \textbf{Topological field theory}: They showed that the geometric Langlands correspondence arises from the duality between two 3d TQFTs.
\item \textbf{Higher categories}: They formulated the correspondence in the language of $(\infty, 2)$-categories, which is the natural setting for categorical Langlands duality.
\item \textbf{Factorization algebras}: They introduced factorization algebras as a bridge between quantum field theory and the geometric Langlands program.
\end{enumerate}

\subsection{Gaitsgory's Program (2010s--2020s)}

Dennis Gaitsgory has been leading a comprehensive program to prove the geometric Langlands correspondence using the theory of \emph{ind-coherent sheaves} and the \emph{geometric Satake equivalence} \cite{gaitsgory2009}. His approach has produced significant partial results.

\subsection{The 2024 Proof (Arinkin, Gaitsgory, Raskin et al.)}

In February 2024, a major breakthrough was achieved by a collaborative team of nine mathematicians: D. Arinkin, D. Beraldo, J. Campbell, L. Chen, J. Faergeman, D. Gaitsgory, K. Lin, S. Raskin, and N. Rozenblyum. In a monumental series of five preprints totaling over 1,000 pages, they proved the unramified geometric Langlands conjecture for any reductive group $G$ over any smooth projective curve $X$ over an algebraically closed field of characteristic 0 \cite{arinkin2024}.

Their proof established the equivalence of derived categories of D-modules and coherent sheaves under the spectral gluing construction, settling the conjecture after nearly four decades of effort.

\begin{highlightbox}
The unramified geometric Langlands conjecture was proved in 2024 by Dennis Gaitsgory and a collaborative team of nine mathematicians, marking a historic milestone in 21st-century geometry and representation theory.
\end{highlightbox}

\subsection{After the Proof}

The unramified geometric Langlands conjecture is settled for any reductive group $G$ over any smooth projective curve in characteristic 0. The 2024 proof by Arinkin, Gaitsgory, Raskin, and collaborators spans over 1,000 pages across five preprints and establishes the equivalence of derived categories between D-modules on the stack of $G$-bundles and coherent sheaves on the stack of ${}^L\!G$-local systems.

The proof introduced powerful new categorical techniques---particularly the theory of ind-coherent sheaves and the spectral gluing construction---that have become tools in their own right. These methods are being applied to problems in derived algebraic geometry, topological field theory, and representation theory well beyond the original Langlands setting. The proof also deepened the connection between the Langlands program and physics, specifically to Chern--Simons theory and other topological field theories.

The full \emph{ramified} geometric Langlands conjecture---where the curve is allowed to have singularities or marked points---remains open. The ramified case is substantially harder because the local data at each ramification point must be incorporated into the global equivalence. Connections to physics, particularly the precise relationship to topological quantum field theories and conformal field theories, are still being developed. Computational aspects---explicitly constructing the functors that realize the Langlands duality---are largely unexplored.

\section{Impact}

\begin{itemize}
\item \textbf{Derived categories}: The geometric Langlands program has driven the development of derived algebraic geometry and the theory of higher categories.
\item \textbf{Mathematical physics}: The connection to topological field theory (Chern--Simons theory, mirror symmetry) has enriched both mathematics and physics.
\item \textbf{Geometric Satake}: The geometric Satake equivalence (Lusztig, Ginzburg, Mirkovic--Vilonen) is a cornerstone result connecting representation theory to geometry \cite{mirkovic2007}.
\item \textbf{Operator algebras}: Connections to vertex operator algebras and conformal field theory (the work of Frenkel, Ben-Zvi, and others) \cite{frenkel2007}.
\end{itemize}

\section{Exercises}

\begin{exercise}[Basic]
What is the Langlands dual group $^LG$ of $G = SL_n$? Of $G = SO_n$? Of $G = Sp_{2n}$?
\end{exercise}

\begin{exercise}[Intermediate]
Explain the Hecke eigensheaf condition in the geometric Langlands program. How does it generalize the eigenform condition in the classical theory?
\end{exercise}

\begin{exercise}[Challenge]
What is the geometric Satake equivalence? How does it relate the representation theory of $^LG$ to the geometry of the affine Grassmannian of $G$? \cite{mirkovic2007}
\end{exercise}

\begin{exercise}
Why are higher categories ($\infty$-categories) the natural language for the geometric Langlands program? What goes wrong with ordinary categories?
\end{exercise}


\begin{exercise}[Computational]
Write a Python/SageMath script to explore a concept from this chapter. See the appendix for starter code.
\end{exercise}

\section{Timeline}

\begin{tabularx}{\linewidth}{lX}
\toprule
\textbf{Year} & \textbf{Event} \\ \midrule
1987 & Drinfeld proposes the geometric Langlands program \cite{drinfeld1987} \\ 
1990s & Beilinson and Drinfeld lay the foundations \cite{beilinson2004} \\ 
1998--2000 & Lusztig, Ginzburg, Mirkovic--Vilonen prove geometric Satake \cite{mirkovic2007} \\
2000s & Laumon, Rothstein, and others develop the theory further \\ 
2010s & Ben-Zvi, Francis, Nadler develop the TQFT approach \cite{benzvi2014} \\ 
2010s--2020s & Gaitsgory's comprehensive program on ind-coherent sheaves \cite{gaitsgory2009} \\ 
2024 & Arinkin, Gaitsgory, Raskin et al. prove the unramified geometric Langlands conjecture \cite{arinkin2024} \\ \bottomrule
\end{tabularx}

\section{Citations}

\begin{enumerate}
\item Arinkin, D., Beraldo, D., Campbell, J., Chen, L., Faergeman, J., Gaitsgory, D., Lin, K., Raskin, S., \& Rozenblyum, N. (2024). ``The unramified geometric Langlands conjecture.'' arXiv:2402.17294.
\item Ben-Zvi, D., Francis, J., \& Nadler, D. (2014). ``Integral transforms and Drinfeld centers in derived algebraic geometry.'' Publications Math\'ematiques de l'IH\'ES, 121(1), 1--137.
\item Gaitsgory, D. (2009). ``Construction of central elements in the affine Hecke algebra via nearby cycles.'' Inventiones Mathematicae, 170(2), 297--347.
\item Mirkovic, I., \& Vilonen, K. (2007). ``Geometric Langlands duality and representations of algebraic groups over commutative rings.'' Annals of Mathematics, 166(1), 95--143.
\item Frenkel, E. (2007). ``Langlands correspondence for loop groups.'' Cambridge Tracts in Mathematics, 173.
\item Beilinson, A., \& Drinfeld, V. (2004). ``Chern--Simons sheaves and geometrical Langlands program.'' In BAMS Surveys.
\item Drinfeld, V. G. (1987). ``Letter to A. B. Volokitin.'' (Published in the collection of Drinfeld's works).
\item Lusztig, G. (1998). ``Cuspidal local systems and graded Hecke algebras.'' Publications Math\'ematiques de l'IH\'ES.
\item Ginzburg, V. (1998). ``Perverse sheaves on a loop group and Langlands duality.'' Preprint.
\end{enumerate}